% essay.tex  —  GENED 1090 Capstone Essay
%
% Compile sequence:
%   pdflatex essay
%   biber essay
%   pdflatex essay
%   pdflatex essay
%
% Requires: TeX Live (or MiKTeX) with biblatex-mla and biber.

\documentclass[12pt]{article}

% ── Font & encoding ───────────────────────────────────────────────────────────
\usepackage[T1]{fontenc}
\usepackage[utf8]{inputenc}
\usepackage{mathptmx}          % Times New Roman for body text and math

% ── Page geometry ─────────────────────────────────────────────────────────────
\usepackage[margin=1in]{geometry}

% ── Line spacing ──────────────────────────────────────────────────────────────
\usepackage{setspace}
\doublespacing

% ── Headers / footers ─────────────────────────────────────────────────────────
% MLA: "Lastname #" top-right on every page, no header rule.
% The name/course/date block is standard MLA body text, not a fancyhdr header.
\usepackage{fancyhdr}
\setlength{\headheight}{14.5pt}
\pagestyle{fancy}
\fancyhf{}
\fancyhead[R]{Ramesh \thepage}
\renewcommand{\headrulewidth}{0pt}

% ── Bibliography (MLA) ────────────────────────────────────────────────────────
\usepackage[style=mla,backend=biber]{biblatex}
\addbibresource{bibliography.bib}

% ── Miscellaneous ─────────────────────────────────────────────────────────────
\usepackage{microtype}          % improved justification
\usepackage{csquotes}           % context-sensitive quotation marks (biblatex)

% ─────────────────────────────────────────────────────────────────────────────

\begin{document}

% MLA info block: body text, top-left of page 1 only.
\noindent Pranav Ramesh\\
GENED 1090\\
05/11/2026

% MLA title: centered, plain — no bold, no italic, no underline.
\begin{center}
  Surface and Residue: \textit{Generative Palimpsest} as Digital Book
\end{center}

\noindent In 1229~\textsc{ce}, an unnamed Byzantine scribe received twelve
leaves of parchment from a monastery library, washed them with milk and oat
bran, and copied over them a Greek Orthodox prayer book, a Euchologion for
use in the Divine Liturgy \parencite{netz2007}. The text he covered had been
written some nine hundred years earlier in the hand of Archimedes. Parchment was costly, and reuse was rational. The writing he erased descended
into the skin, held by iron gall ink that had bitten too deep for any sponge
to remove, and waited. In 1906 the Danish
classicist J.~L. Heiberg recognized the prayer book as a palimpsest containing
works of Archimedes. In 2001, multispectral imaging at the Walters Art Museum
in Baltimore recovered the proof that Archimedes had worked with actual
infinity, a concept not formally rediscovered until Newton and Leibniz,
fourteen centuries later.

\textit{Generative Palimpsest} takes this story as its governing metaphor and
structural model. It is a web-based interactive book, rendered on a laptop or
tablet screen as a two-page spread, with six pages of layered text across
chapters on Memory, Erasure, the Archimedes Palimpsest, Preservation, Loss,
and the Generative Scribe. A parchment-textured veil covers each page; the
reader scrapes it away with the cursor, revealing text beneath. When a cleared
zone exceeds sixty percent transparency, an artificial intelligence language
model surfaces a new prose fragment in a distinct color, text generated in the
moment of the reader's looking. \textit{Generative Palimpsest}
deploys the digital screen as a writing platform that behaves structurally like
the reused medieval parchment page: a surface written over many times, from
which older layers emerge when the new is scraped away. In enacting this
structure, the book makes a claim about large language models: they are shaped
by their training data and surface residue from it the way a palimpsest
surfaces ink.

The platform of \textit{Generative Palimpsest} is the digital screen; its
writing support is the HTML5 Canvas element, rendered in a web browser with a
parchment texture and an ink palette derived from aged manuscripts. In the
taxonomy Amaranth Borsuk develops in \textit{The Book}, this places the project
within a long history of material experimentation in which each writing platform
brings distinct affordances to the act of reading. Borsuk traces a lineage from
wax tablet to papyrus scroll to parchment codex to paper to digital screen,
arguing that the book has always been a site of experimentation and
transformation, never settling into a fixed form \parencite{borsuk2018}. Only a
computational surface can hold generative layers, text that the reader's act
of looking calls into being. Clay, papyrus, parchment, and paper are fixed
surfaces. The screen is the first writing platform that can write back.

The scraping mechanic, implemented through the Canvas \texttt{destination-out}
compositing operation, which erases the opaque veil where the cursor drags, is
the computational equivalent of the scribe's wet sponge or scraping knife.
Richard H.~Rouse and Mary A.~Rouse document the dense economic network of
parchment sellers, scribes, illuminators, and stationers that sustained
commercial book production in medieval Paris, a trade in material as much as
in text, in which parchment was costly and reuse was rational \parencite{rouse2000}.
The gesture of
scraping has been translated from vellum to the digital interface; the
relationship between reader and surface, between concealment and disclosure,
remains.

Each of the six pages contains three pre-composed text layers at different
opacities and rotations, mimicking the palimpsest's characteristic overwriting.
The layers are drawn from heterogeneous sources: the Latin Vulgate Psalms
(Psalm~89, \textit{c.}~390~\textsc{ce}), Sappho's Fragment~16 in Anne Carson's
translation, a verse from the Qur\textquoteleft an (3:185), scholarly accounts
of the Archimedes Palimpsest, archival descriptions of library fires, and my own prose. Each non-English source layer (the Vulgate Latin and the
Arabic Qur\textquoteleft anic verse) is accompanied by an English translation
in the footnote strip at the base of its page; the Sappho fragment is presented
exclusively in Carson's English rendering. Even before scraping begins, all layers are partially visible through the semi-transparent veil.

The codicological term for what lies beneath is \textit{scriptio inferior}, the
lower writing; what is visible on the surface is the \textit{scriptio superior}.
Both are present simultaneously, at different opacities, before any interaction
begins. The choice of four distinct serif typefaces (IM Fell English, Spectral,
Lora, and Cormorant Garamond, each assigned to a different layer) mimics the
way different scribal hands are distinguishable by paleographers examining a
manuscript. M.~B. Parkes, in his studies of medieval scribal practice, observes
that layout and decoration ``function like punctuation'': they are part of the
presentation of a text that both facilitates its use by a reader and interprets
the text as transmitted to the scribe
\parencite{parkes1991}. The typographic assignment in \textit{Generative Palimpsest} operates accordingly:
font becomes a codicological cue, distinguishing the eleventh-century liturgical
stratum from the modern-archival stratum from my own voice.

The act of scraping makes the reader a conservator. The cursor performs what
ultraviolet light and multispectral imaging perform on a material palimpsest: it
renders the surface transparent, disclosing what is beneath. Johanna Drucker,
analyzing what the electronic book must inherit from the codex, argues that the
traditional manuscript is ``already, in an important and suggestive way,
virtual'': its apparatus of marginalia, running titles, catchwords, and quire
signatures constitutes a navigational system for moving through space not
literally present on any single page \parencite{drucker2008}. The book applies
this observation
literally. It includes foliation (fol.~1r through fol.~6r), three quire
signatures (\textit{a~i}, \textit{b~i}, \textit{c~i}), catchwords at the foot
of each page anticipating the text of the next, running titles in the top
margin, outer marginalia, hidden marginalia revealed only by scraping the lower
margin, and manicule reader-marks concealed beneath the veil. These are the organizing apparatus of the medieval manuscript
translated into a computational environment where they remain meaningfully
operative. The catchwords guide the eye across the virtual spine, the quire
signatures mark the gathering structure, the marginalia comment on the layers
they flank. A UV mode, toggled by a button in the interface, reveals all three
text layers simultaneously in distinct luminous hues, simulating the
multispectral imaging used to recover the Archimedes undertext at the Walters.
What the conservators accomplished with physics, the reader accomplishes with a
keystroke.

When a cleared zone exceeds sixty percent transparency, the application calls a
large language model that generates a short prose fragment instructed to surface
``as a recovered fragment, something that was nearly lost and not entirely
coherent.'' This fourth layer reconstitutes language from the residue of its training,
assembling fragments rather than retrieving stored passages. Parkes
describes how a medieval scribe's errors became, in subsequent copies, new
readings, ``and sometimes, new meanings,'' because the scribe was a conduit for
texts he often did not fully understand, in languages he may have known
imperfectly, his attention wandering, his hand tired \parencite{parkes1991}.
The scribe introduced
variants. So does the model: it echoes, distorts, or quietly contradicts what
is visible on the page, generating language that is original in the act of
generation, borrowed from nothing prior. Both the scribe and the model are machines
for the transmission of language that distort what they transmit. What we call a
variant in the scribe, we call a hallucination in the model; the mechanism
differs, but both occupy the same structural position in the transmission of
language.

As Netz and Noel recount, the prayer book that Heiberg identified concealed
seven treatises by Archimedes, including \textit{The Method of Mechanical
  Theorems}, the only surviving witness to Archimedes' engagement with infinitely
large sets, and his closest approach to integral calculus. The manuscript was
sold at Christie's in 1998, deposited at the Walters Art Museum in Baltimore,
and subjected to multispectral imaging that recovered approximately eighty
percent of the undertext; the remaining twenty percent required synchrotron
X-ray fluorescence imaging to map the iron atoms in Archimedes' ink. Some
lacunae remain permanent. This is the third chapter of \textit{Generative
  Palimpsest}: the Archimedes story is the book's central case study, the
palimpsest in which what descended into the skin changed the history of
mathematics \parencite{netz2007}.

L.~D. Reynolds and N.~G. Wilson, tracing the transmission of Greek and Latin
literature through the medieval period, observe that of the seventy plays of
Aeschylus, seven survive; of the one hundred and twenty-three of Sophocles,
seven; of the ninety of Euripides, eighteen, and we have more of Euripides
because the Byzantine schoolmasters chose him for study \parencite{reynolds1991}.
What was chosen was copied, and what was copied survived. The fifth chapter of \textit{Generative
  Palimpsest}, on Loss, opens with a fragment from the Oxyrhynchus papyrus
P.Oxy.~1787, a partially legible Sappho poem in which lacunae are marked with
brackets, signaling the absence of papyrus that is no longer there. Anne
Carson, in \textit{Nox}, describes reading such a fragmentary text as ``reading
in a very dark room'': knowing ``there is furniture,'' putting your hands out
and touching ``the cool edge of a table, the rough weave of a chair,''
constructing the room from the parts you can reach, ``never certain the room is
the shape you imagine'' \parencite{carson2010}. The model's training corpus is also a curriculum,
assembled by engineers making decisions about what text to include, decisions as
consequential for what will be remembered as any medieval selection, and less
transparent in their criteria. The brackets in the
papyrological transcription and the model's silences mark the same kind of
absence.

\textit{Generative Palimpsest} is a philosophical book whose argument is enacted
through its structure: the reader grasps the claim by scraping, before any
sentence names it. N.~Katherine Hayles, in \textit{Writing Machines}, develops the
concept of the technotext, a literary work that foregrounds the inscription
technology that produces it, making the relationship between medium and content
legible as argument in itself \parencite{hayles2002}. \textit{Generative
  Palimpsest} is a technotext
in this sense. It is a palimpsest that generates, asking the reader to perform, physically,
the analogy between digital scraping and manuscript conservation. Three of the eighteen text layers are my
own prose, attributed to ``Pranav Ramesh, 2025,'' sitting at the same opacity
as the Vulgate Psalm and the Sappho fragment, available only to the reader who
scrapes to find them.
My authorship is submerged in the same way: my prose sits beneath the veil
as Archimedes sits beneath the Euchologion, discoverable only by a reader who
scrapes. The colophon on the final page reads: ``This book has
no final layer. It will continue to generate as long as you continue to read.''
The \textit{explicit}, the medieval manuscript's closing formula marking
completion, is deferred.

In its material apparatus, \textit{Generative Palimpsest} reproduces the full
parchment codex: a two-page spread across three gatherings, each opening with a
quire signature; foliation in the style of recto leaves; catchwords, running
titles, marginalia, footnote strips, a table of contents, and a colophon. The
page-turn is animated via clip-path transition in the manner of a physical
opening; the spine is rendered as a gradient gutter between pages. A
paleographer would recognize the organizational apparatus at a glance. Every
reader who opens \textit{Generative Palimpsest} receives a different book, because the AI-generated fourth layer is new on each
visit, constructed in response to what that particular reader has uncovered on
that particular page. The medieval codex preserved its \textit{ordinatio}, the
overall arrangement of text and paratext, across many copies, each introducing
small scribal variants but preserving the structure of the exemplar
\parencite{parkes1991}. \textit{Generative Palimpsest} has no exemplar. Every
copy is unique in the way every palimpsest is unique: the residue that survives
depends on the particular conjunction of ink, vellum, sponge, and scribal hand.

Borsuk argues that the history of the book records our ``highly contingent
cultural ideals of authorship and art,'' a record of ongoing material
transformation in which each new platform reconfigures what reading and writing
can mean \parencite{borsuk2018}. The imagined reader is someone who has held a medieval manuscript and
used a language model, someone who can feel, in the act of scraping, the
equivalence of the scholar's UV lamp and the reader's cursor: both press
through a visible surface to what lies beneath, and find something made in
the act of looking. The palimpsest teaches that erasure is incomplete; ink descends into skin and
waits. \textit{Generative
  Palimpsest} extends this lesson to the present moment and asks whether what a
language model surfaces is better understood as retrieval or as residue. The
book declines to answer directly; it asks the reader to scrape.

% Ensure all 8 sources appear in Works Cited regardless of citation style.
\nocite{netz2007,parkes1991,reynolds1991,borsuk2018,hayles2002,drucker2008,rouse2000,carson2010}

\newpage
\printbibliography[title={Works Cited}]

\end{document}
